% !TEX TS-program = pdflatex
% Document LaTeX — Modèle de classes backend

\documentclass[11pt,a4paper]{article}

\usepackage[T1]{fontenc}
\usepackage[utf8]{inputenc}
\usepackage[french]{babel}
\usepackage{geometry}
\usepackage{hyperref}
\usepackage{enumitem}
\usepackage{tabularx}
\usepackage{longtable}
\usepackage{booktabs}
\usepackage{xcolor}

\geometry{margin=2.5cm}
\hypersetup{
  colorlinks=true,
  linkcolor=black,
  urlcolor=blue,
  citecolor=black
}

\setlist[itemize]{noitemsep, topsep=4pt}
\setlist[description]{style=nextline, noitemsep, topsep=4pt}

\title{Projet transversal (MIF10) — Groupe 15\\\vspace{0.2cm}Modèle de classes du backend}
\author{Groupe 15}
\date{18 mars 2026}

\begin{document}
\maketitle

\section{Contexte et objectifs}
Le projet vise la réalisation d’une application web de gestion d’évènements participatifs (type \textit{Meetup}) centrée sur des évènements caritatifs autour du cinéma.

Les fonctionnalités attendues côté backend incluent :
\begin{itemize}
  \item consultation, recherche et filtrage des évènements (visiteurs) ;
  \item création de compte, authentification, favoris, inscriptions et commentaires (utilisateurs connectés) ;
  \item demande de création d’association (utilisateurs connectés) ;
  \item création, publication, modification, annulation d’évènements et gestion d’illustrations (utilisateurs rattachés à une association) ;
  \item consultation de la liste des inscrits à un évènement (association).
\end{itemize}

\section{Hypothèses et décisions de modélisation}
\begin{itemize}
  \item \textbf{Pas de notion d’admin dédiée} : un utilisateur « éditeur » est un utilisateur rattaché à une \texttt{Association}.
  \item La publication est gérée par un \textbf{statut} \texttt{DRAFT/PUBLISHED/CANCELLED}.
  \item L’annulation ne déclenche pas de notification automatique aux participants (simple changement de statut visible côté UI).
  \item Un utilisateur peut avoir \textbf{au plus un commentaire par évènement} (cohérent avec « modifier mon commentaire »).
  \item Un utilisateur est rattaché à \textbf{au plus une association} (cardinalité 0..1).
  \item Les utilisateurs peuvent faire une \textbf{demande de création d’association} (workflow \texttt{PENDING/APPROVED/REJECTED}).
  \item Les évènements ont des \textbf{catégories hiérarchiques} (catégorie parent / catégories enfants).
\end{itemize}

\section{Glossaire (types et conventions)}
\begin{description}
  \item[\texttt{UUID}] Identifiant unique (stockage DB : \texttt{uuid} ou \texttt{char/varchar}).
  \item[\texttt{datetime}] Date/heure (stockage DB : \texttt{timestamp} avec fuseau conseillé).
  \item[\texttt{text}] Chaîne potentiellement longue.
\end{description}

\section{Énumérations et objets valeur}
\subsection{\texttt{EventStatus}}
\begin{itemize}
  \item \texttt{DRAFT} : évènement non visible publiquement.
  \item \texttt{PUBLISHED} : évènement visible publiquement.
  \item \texttt{CANCELLED} : évènement visible publiquement, marqué comme annulé.
\end{itemize}

\subsection{\texttt{AssociationRequestStatus}}
\begin{itemize}
  \item \texttt{PENDING} : demande en attente de traitement.
  \item \texttt{APPROVED} : demande acceptée (création d’une association).
  \item \texttt{REJECTED} : demande refusée (motif optionnel).
\end{itemize}

\subsection{\texttt{Location} (objet valeur)}
La localisation d’un évènement est modélisée comme un objet valeur (ou des colonnes directement dans \texttt{Event}).
Champs suggérés : \texttt{name}, \texttt{addressLine1}, \texttt{postalCode}, \texttt{city}, \texttt{country}, \texttt{lat}, \texttt{lng} (optionnels).

\section{Modèle de domaine (vue d’ensemble)}
\subsection{Entités}
\begin{itemize}
  \item \texttt{User} : compte utilisateur.
  \item \texttt{Association} : organisation qui publie les évènements.
  \item \texttt{AssociationRequest} : demande de création d’association.
  \item \texttt{Event} : évènement caritatif (projection/activité autour d’un film).
  \item \texttt{EventImage} : image d’illustration d’un évènement.
  \item \texttt{Favorite} : favori d’un utilisateur sur un évènement.
  \item \texttt{Signup} : participation/inscription d’un utilisateur à un évènement.
  \item \texttt{Comment} : commentaire d’un utilisateur sur un évènement.
  \item \texttt{EventCategory} : catégorie (hiérarchique).
  \item \texttt{EventCategoryLink} : association évènement--catégorie.
\end{itemize}

\subsection{Relations principales}
\begin{itemize}
  \item Une \texttt{Association} possède plusieurs \texttt{User} (membres/éditeurs).
  \item Une \texttt{Association} publie plusieurs \texttt{Event}.
  \item Un \texttt{Event} possède plusieurs \texttt{EventImage}.
  \item \texttt{Favorite}, \texttt{Signup} et \texttt{Comment} sont des relations entre \texttt{User} et \texttt{Event}.
  \item Les catégories sont hiérarchiques : une \texttt{EventCategory} peut avoir un parent et plusieurs enfants.
\end{itemize}

\section{Spécification des classes / tables}
Les sections suivantes décrivent les champs, les contraintes et les règles de gestion.

\subsection{\texttt{User}}
\textbf{Rôle métier} : utilisateur connecté (peut liker, s’inscrire, commenter). S’il est rattaché à une association, il peut également gérer des évènements.

\begin{longtable}{@{}p{0.28\textwidth}p{0.18\textwidth}p{0.46\textwidth}@{}}
\toprule
\textbf{Champ} & \textbf{Type} & \textbf{Contraintes / remarques}\\
\midrule
\endhead
\texttt{id} & UUID & Clé primaire.\\
\texttt{email} & string & Unique, format email.\\
\texttt{passwordHash} & string & Hash (jamais le mot de passe en clair).\\
\texttt{displayName} & string & Nom affiché.\\
\texttt{hasAssociation} & boolean & Vrai si l’utilisateur est rattaché à une association.\\
\texttt{associationId} & UUID? & Nullable ; non-null si \texttt{hasAssociation=true}.\\
\texttt{createdAt} & datetime & Audit.\\
\texttt{updatedAt} & datetime & Audit.\\
\texttt{disabledAt} & datetime? & Optionnel (désactivation).\\
\bottomrule
\end{longtable}

\textbf{Contraintes clés}
\begin{itemize}
  \item \texttt{email} unique.
  \item Cohérence : \texttt{hasAssociation=true} $\Rightarrow$ \texttt{associationId} non-null.
\end{itemize}

\subsection{\texttt{Association}}
\textbf{Rôle métier} : entité porteuse des évènements (publication et promotion).

\begin{longtable}{@{}p{0.28\textwidth}p{0.18\textwidth}p{0.46\textwidth}@{}}
\toprule
\textbf{Champ} & \textbf{Type} & \textbf{Contraintes / remarques}\\
\midrule
\endhead
\texttt{id} & UUID & Clé primaire.\\
\texttt{name} & string & Nom officiel.\\
\texttt{description} & text? & Optionnel.\\
\texttt{websiteUrl} & string? & Optionnel.\\
\texttt{logoUrl} & string? & Optionnel.\\
\texttt{createdAt} & datetime & Audit.\\
\texttt{updatedAt} & datetime & Audit.\\
\bottomrule
\end{longtable}

\subsection{\texttt{AssociationRequest}}
\textbf{Rôle métier} : demande d’un utilisateur pour créer une association.

\begin{longtable}{@{}p{0.28\textwidth}p{0.18\textwidth}p{0.46\textwidth}@{}}
\toprule
\textbf{Champ} & \textbf{Type} & \textbf{Contraintes / remarques}\\
\midrule
\endhead
\texttt{id} & UUID & Clé primaire.\\
\texttt{requestedByUserId} & UUID & FK vers \texttt{User}.\\
\texttt{name} & string & Nom proposé.\\
\texttt{description} & text? & Optionnel.\\
\texttt{websiteUrl} & string? & Optionnel.\\
\texttt{status} & AssociationRequestStatus & \texttt{PENDING|APPROVED|REJECTED}.\\
\texttt{reviewedAt} & datetime? & Optionnel.\\
\texttt{reviewNote} & string? & Optionnel (motif).\\
\texttt{createdAt} & datetime & Audit.\\
\texttt{updatedAt} & datetime & Audit.\\
\bottomrule
\end{longtable}

\textbf{Contraintes}
\begin{itemize}
  \item Recommandé : un utilisateur ne peut avoir qu’une demande \texttt{PENDING} à la fois.
  \item À l’approbation : création d’une \texttt{Association} et rattachement du \texttt{User} demandeur.
\end{itemize}

\subsection{\texttt{Event}}
\textbf{Rôle métier} : évènement caritatif autour d’un film.

\begin{longtable}{@{}p{0.28\textwidth}p{0.18\textwidth}p{0.46\textwidth}@{}}
\toprule
\textbf{Champ} & \textbf{Type} & \textbf{Contraintes / remarques}\\
\midrule
\endhead
\texttt{id} & UUID & Clé primaire.\\
\texttt{title} & string & Titre affiché de l’évènement.\\
\texttt{movieTitle} & string & Supporte la recherche « nom du film ».\\
\texttt{description} & text & Description détaillée.\\
\texttt{startsAt} & datetime & Date/heure de début.\\
\texttt{endsAt} & datetime? & Optionnel.\\
\texttt{location} & Location & Objet valeur (ou colonnes).\\
\texttt{capacity} & int? & Optionnel ; si null, capacité illimitée.\\
\texttt{status} & EventStatus & \texttt{DRAFT|PUBLISHED|CANCELLED}.\\
\texttt{cancelReason} & string? & Optionnel ; utile si annulé.\\
\texttt{associationId} & UUID & Association qui publie l’évènement.\\
\texttt{createdByUserId} & UUID & Doit référencer un utilisateur rattaché à une association.\\
\texttt{publishedAt} & datetime? & Défini lors de la publication.\\
\texttt{createdAt} & datetime & Audit.\\
\texttt{updatedAt} & datetime & Audit.\\
\bottomrule
\end{longtable}

\textbf{Règles métier}
\begin{itemize}
  \item Visibilité publique : \texttt{status \in \{PUBLISHED, CANCELLED\}}.
  \item Publication : \texttt{DRAFT \to PUBLISHED} (renseigne \texttt{publishedAt}).
  \item Annulation : \texttt{PUBLISHED \to CANCELLED} (optionnellement renseigne \texttt{cancelReason}).
  \item Autorisation : seuls les utilisateurs rattachés à une association peuvent créer/publier/modifier/annuler.
  \item Capacité : si \texttt{capacity} est définie, le nombre de \texttt{Signup} ne doit pas dépasser \texttt{capacity}.
\end{itemize}

\subsection{\texttt{EventImage}}
\textbf{Rôle métier} : illustration d’un évènement.

\begin{longtable}{@{}p{0.28\textwidth}p{0.18\textwidth}p{0.46\textwidth}@{}}
\toprule
\textbf{Champ} & \textbf{Type} & \textbf{Contraintes / remarques}\\
\midrule
\endhead
\texttt{id} & UUID & Clé primaire.\\
\texttt{eventId} & UUID & FK vers \texttt{Event}.\\
\texttt{url} & string & URL publique ou clé de stockage.\\
\texttt{altText} & string? & Optionnel.\\
\texttt{position} & int & Ordre d’affichage.\\
\texttt{createdAt} & datetime & Audit.\\
\bottomrule
\end{longtable}

\textbf{Contraintes}
\begin{itemize}
  \item \texttt{(eventId, position)} unique (recommandé).
\end{itemize}

\subsection{\texttt{Favorite}}
\textbf{Rôle métier} : favori d’un utilisateur sur un évènement.

\begin{longtable}{@{}p{0.28\textwidth}p{0.18\textwidth}p{0.46\textwidth}@{}}
\toprule
\textbf{Champ} & \textbf{Type} & \textbf{Contraintes / remarques}\\
\midrule
\endhead
\texttt{id} & UUID & Clé primaire.\\
\texttt{userId} & UUID & FK vers \texttt{User}.\\
\texttt{eventId} & UUID & FK vers \texttt{Event}.\\
\texttt{createdAt} & datetime & Audit.\\
\bottomrule
\end{longtable}

\textbf{Contraintes}
\begin{itemize}
  \item Unicité : \texttt{(userId, eventId)} unique.
  \item Autorisé uniquement pour un utilisateur authentifié.
\end{itemize}

\subsection{\texttt{Signup}}
\textbf{Rôle métier} : inscription/participation à un évènement.

\begin{longtable}{@{}p{0.28\textwidth}p{0.18\textwidth}p{0.46\textwidth}@{}}
\toprule
\textbf{Champ} & \textbf{Type} & \textbf{Contraintes / remarques}\\
\midrule
\endhead
\texttt{id} & UUID & Clé primaire.\\
\texttt{userId} & UUID & FK vers \texttt{User}.\\
\texttt{eventId} & UUID & FK vers \texttt{Event}.\\
\texttt{createdAt} & datetime & Audit.\\
\bottomrule
\end{longtable}

\textbf{Contraintes}
\begin{itemize}
  \item Unicité : \texttt{(userId, eventId)} unique.
  \item Autorisé uniquement si \texttt{Event.status = PUBLISHED}.
  \item Si \texttt{capacity} est définie, refuser l’inscription lorsque l’évènement est complet.
\end{itemize}

\subsection{\texttt{Comment}}
\textbf{Rôle métier} : commentaire d’un utilisateur sur un évènement.

\begin{longtable}{@{}p{0.28\textwidth}p{0.18\textwidth}p{0.46\textwidth}@{}}
\toprule
\textbf{Champ} & \textbf{Type} & \textbf{Contraintes / remarques}\\
\midrule
\endhead
\texttt{id} & UUID & Clé primaire.\\
\texttt{userId} & UUID & FK vers \texttt{User}.\\
\texttt{eventId} & UUID & FK vers \texttt{Event}.\\
\texttt{body} & text & Contenu du commentaire.\\
\texttt{createdAt} & datetime & Audit.\\
\texttt{updatedAt} & datetime & Audit.\\
\texttt{deletedAt} & datetime? & Optionnel (soft delete).\\
\bottomrule
\end{longtable}

\textbf{Contraintes}
\begin{itemize}
  \item Recommandé : \texttt{(userId, eventId)} unique (un commentaire par évènement et par utilisateur).
  \item Modification/suppression : par l’auteur (et optionnellement par un utilisateur d’association si vous décidez une modération).
\end{itemize}

\subsection{\texttt{EventCategory}}
\textbf{Rôle métier} : classification des évènements, sous forme hiérarchique (parent/enfants).

\begin{longtable}{@{}p{0.28\textwidth}p{0.18\textwidth}p{0.46\textwidth}@{}}
\toprule
\textbf{Champ} & \textbf{Type} & \textbf{Contraintes / remarques}\\
\midrule
\endhead
\texttt{id} & UUID & Clé primaire.\\
\texttt{name} & string & Nom lisible.\\
\texttt{slug} & string & Unique, pratique pour URLs.\\
\texttt{parentId} & UUID? & FK vers \texttt{EventCategory}, nullable.\\
\texttt{createdAt} & datetime & Audit.\\
\texttt{updatedAt} & datetime & Audit.\\
\bottomrule
\end{longtable}

\textbf{Contraintes}
\begin{itemize}
  \item \texttt{parentId} nullable pour les catégories racines.
  \item Éviter les cycles (contrainte applicative).
\end{itemize}

\subsection{\texttt{EventCategoryLink}}
\textbf{Rôle métier} : associer un évènement à une ou plusieurs catégories.

\begin{longtable}{@{}p{0.28\textwidth}p{0.18\textwidth}p{0.46\textwidth}@{}}
\toprule
\textbf{Champ} & \textbf{Type} & \textbf{Contraintes / remarques}\\
\midrule
\endhead
\texttt{id} & UUID & Clé primaire.\\
\texttt{eventId} & UUID & FK vers \texttt{Event}.\\
\texttt{categoryId} & UUID & FK vers \texttt{EventCategory}.\\
\texttt{createdAt} & datetime & Audit.\\
\bottomrule
\end{longtable}

\textbf{Contraintes}
\begin{itemize}
  \item Unicité : \texttt{(eventId, categoryId)} unique.
  \item Option : imposer une seule catégorie par évènement via contrainte unique sur \texttt{eventId}.
\end{itemize}

\section{Vues / requêtes utiles (implémentation)}
\subsection{Dashboard utilisateur}
\begin{itemize}
  \item \textbf{Favoris} : événements joints avec \texttt{Favorite} filtrés par \texttt{userId} (exclure les \texttt{DRAFT}).
  \item \textbf{Participations} : événements joints avec \texttt{Signup} filtrés par \texttt{userId}.
\end{itemize}

\subsection{Liste des inscrits à un évènement (association)}
\begin{itemize}
  \item Pour un \texttt{eventId}, lister les \texttt{User} joints avec \texttt{Signup}.
\end{itemize}

\section{Recherche et filtrage (visiteur)}
\subsection{Critères minimum}
\begin{itemize}
  \item Recherche « film » : sur \texttt{Event.movieTitle}.
  \item Filtre date : intervalle \texttt{[from, to]} sur \texttt{Event.startsAt}.
  \item Filtre lieu : au minimum \texttt{city} (option : code postal).
  \item Filtre catégorie : par catégorie (option : inclure les sous-catégories).
  \item Visibilité : \texttt{status \in \{PUBLISHED, CANCELLED\}}.
\end{itemize}

\subsection{Index recommandés (base de données)}
\begin{itemize}
  \item \texttt{Event(status, startsAt)}
  \item \texttt{Event(movieTitle)} (ou index full-text)
  \item \texttt{Event(city)} si la ville est stockée en colonne
  \item \texttt{Favorite(userId, eventId)} unique
  \item \texttt{Signup(userId, eventId)} unique
  \item \texttt{Comment(userId, eventId)} unique (si contrainte activée)
  \item \texttt{EventCategory(slug)} unique
  \item \texttt{EventCategory(parentId)}
  \item \texttt{EventCategoryLink(eventId, categoryId)} unique
  \item \texttt{AssociationRequest(requestedByUserId, status)}
\end{itemize}

\end{document}
